\documentclass[11pt]{article}
% Shared print layout. The web's provisional design and MathJax are independent.
\usepackage[utf8]{inputenc}
\usepackage[T1]{fontenc}
\usepackage[spanish,es-nodecimaldot]{babel}
\usepackage{lmodern}
\usepackage{amsmath,amssymb}
\usepackage{graphicx}
\usepackage{geometry}
\geometry{a4paper,left=20mm,right=20mm,top=25mm,bottom=22mm,headheight=16pt,headsep=8mm}
\usepackage{fancyhdr,lastpage}
\usepackage{enumitem}
\usepackage{xcolor}
\usepackage{pgfplots}
\pgfplotsset{compat=1.18}
\usepgfplotslibrary{fillbetween}
\usepackage[hidelinks,pdfencoding=auto]{hyperref}
\hypersetup{pdflang=es-ES,pdfauthor={},pdfsubject={Matemáticas Aplicadas a las Ciencias Sociales II}}
\pagestyle{fancy}
\fancyhf{}
\fancyhead[L]{\small Matemáticas Aplicadas a las Ciencias Sociales II}
\fancyhead[R]{\small 2.\textsuperscript{o} Bachillerato}
\fancyfoot[L]{\footnotesize Distribuciones I | Edición revisada 1.0}
\fancyfoot[R]{\footnotesize Página \thepage\ de \pageref{LastPage}}
\renewcommand{\headrulewidth}{0.4pt}
\setlength{\parindent}{0pt}
\setlength{\parskip}{0.5em}
\setlist{itemsep=0.3em,topsep=0.3em}
\newcommand{\apartado}[1]{\subsection*{#1}}
\pgfplotsset{
  revision/.style={
    width=0.94\linewidth,height=7.5cm,
    axis lines=middle,grid=major,grid style={gray!18},
    tick label style={font=\small},label style={font=\small},
    xlabel={$x$},ylabel={$y$},samples=160,
    legend style={at={(0.5,-0.16)},anchor=north,draw=none,font=\small},
    clip=true
  }
}

\hypersetup{pdftitle={Distribuciones I - Soluciones revisadas}}
\begin{document}
\section*{Ejercicio 1: baloncesto y binomial}
Soluciones revisadas a partir de las incluidas en F003. Se conserva el modelo binomial del original, con independencia entre ensayos y probabilidad constante. El tamaño de una población finita no está especificado.
\[X\sim B(6,0{,}25),\qquad P(X=k)=\binom6k(0{,}25)^k(0{,}75)^{6-k}.\]
\subsection*{(a) Ninguno practica baloncesto}
\[P(X=0)=(0{,}75)^6=\boxed{0{,}177978515625}\approx0{,}1780.\]
\subsection*{(b) Al menos cinco}
\begin{align*}
P(X=5)&=6(0{,}25)^5(0{,}75)=0{,}00439453125,\\
P(X=6)&=(0{,}25)^6=0{,}000244140625.
\end{align*}
\[\boxed{P(X\geq5)=0{,}004638671875}.\]
El original suma los términos ya redondeados y da 0,004644. Se evita ese redondeo intermedio: a cuatro cifras significativas, el resultado es 0,004639.
\subsection*{(c) Al menos uno}
\[P(X\geq1)=1-(0{,}75)^6=\boxed{0{,}822021484375}.\]
A cuatro decimales, 0,8220. Se mantiene la relación exacta de complemento con el apartado (a).
\clearpage
\section*{Ejercicio 2: normal y umbral}
Se interpreta $N(85,20)$ como en los cálculos originales: media 85 y desviación típica 20. Sea $Z=(X-85)/20$ y $\Phi$ la distribución acumulada de $N(0,1)$.
\subsection*{(a) Entre 80 y 120}
\[P(80\leq X\leq120)=\Phi(1{,}75)-\Phi(-0{,}25)\approx0{,}55864717.\]
El porcentaje es $\boxed{55{,}864717\%}\approx55{,}86\%$. La tabla original da 0,5586, coherente con su precisión.
\subsection*{(b) Menos de 75}
\[P(X<75)=\Phi(-0{,}5)\approx\boxed{0{,}30853754}\approx0{,}3085.\]
\subsection*{(c) Bonificación}
Si el 80,23\% inferior no la recibe, el umbral satisface
\[P(X\leq x)=0{,}8023.\]
Con cálculo sin redondeo intermedio,
\[z=\Phi^{-1}(0{,}8023)\approx0{,}84986526,\qquad
\boxed{x\approx101{,}99731\text{ puntos}}.\]
La tabla redondea $z$ a 0,85. Con esa aproximación,
\[x\approx85+20\cdot0{,}85=85+17=\boxed{102\text{ puntos}}.\]
\textbf{Corrección M03:} se conserva el resultado aproximado 102, pero se corrige la igualdad del original que mezcla 0,85, 0,8 y 16,8. El producto correcto es $20\cdot0{,}85=17$; todas las igualdades del desarrollo deben usar el mismo valor.
\clearpage
\section*{Ejercicio 3: proyectos rentables}
Se conserva el modelo de resultados independientes con probabilidad constante:
\[X\sim B(50,0{,}5),\quad\mu=25,\quad\sigma=\sqrt{12{,}5}.\]
El objetivo incluye el caso de exactamente 30 proyectos rentables.
\subsection*{Probabilidad binomial exacta}
\[P(X\geq30)=\sum_{k=30}^{50}\binom{50}{k}(0{,}5)^{50}
\approx\boxed{0{,}10131938}.\]
Es aproximadamente el 10,1319\%.
\subsection*{Aproximación normal del original}
Aquí $np=n(1-p)=25$. Para la aproximación se introduce $Y\sim N(25,\sqrt{12{,}5})$; no se sustituye la distribución exacta de $X$ por una igualdad de modelos.

Con \textbf{corrección de continuidad},
\begin{align*}
P(X\geq30)&\approx P(Y\geq29{,}5),\\
z&=\frac{29{,}5-25}{\sqrt{12{,}5}}\approx1{,}27279221,\\
P(Y\geq29{,}5)&\approx\boxed{0{,}10154589}.
\end{align*}
La tabla del original redondea $z$ a 1,27 y obtiene 0,1020. Es una aproximación de tabla, no el valor exacto binomial. Se conserva esa referencia identificando su precisión.

La diferencia entre normal y binomial, antes de redondear, es aproximadamente 0,00022652. El modelo y los resultados son los del ejercicio, sin introducir criterios de inversión externos.
\end{document}
